\section{Web page classification}\label{sec:webpage_classification}

Web page classification is the process of categorizing web pages in one or more predefined classes based on their content, structure, and visual properties.
This can be considered to be a pre-processing task of information extraction/retrieval, as the method and types of information extraction/retrieval can be different for different types of web pages.
For example, the method of information extraction for a e-commerce website's product page may be different from the method of information extraction for it's contact details page.
Web page classification, while fundamentally can be considered to be a text document classification task, is a unique task due to the complex structure of web pages and the presence of both content and presentation information.\cite{ajose2020systematic}

Web page classification techniques can be broadly categorized into three categories:
\begin{enumerate}
    \item Keyword-based techniques: These techniques rely on the presence of specific keywords or phrases in the web page and its frequency to classify it into a particular category.
    \item Supervised learning techniques: These techniques use labeled training data to train a machine learning model to classify web pages into different categories based on their content and structure.
\end{enumerate}

\subsection{Keyword-based techniques}\label{subsubsec:keyword-based-techniques}
Keyword-based techniques are fairly simple, rule based approaches which do not require a labeled training or training.
Documents are assigned to categories based on the highest match count of specific keywords, or on keyword density (the ratio of the number of occurrences of a keyword to the total number of words in the document).
While such techniques are simple and easy to implement, they can be inaccurate and perform poorly on more complex pages\cite{Procycons_2025}.
More advanced techniques such as Term Frequency - Inverse Document Frequency (TF-IDF) can be used to improve the performance of keyword-based techniques by giving more weight to keywords that are more relevant to a particular category and less weight to common keywords that appear in many categories.
Combining TF-IDF with traditional machine learning classifiers have shows to have strong performance in web page classification tasks\cite{de2019towards}.

\subsection{Supervised learning techniques}\label{subsubsec:supervised-learning-techniques}
Supervised learning techniques for web page classification involve training a machine learning model on a labeled dataset of web pages, where each web page is associated with a specific category or class.
The model learns to classify web pages based on the features extracted from the content and structure of the web pages.
Classical algorithms such as Naive Bayes, Support Vector Machines (SVM), k-Nearest Neighbour (kNN) and Decision Trees have been used for web page classification, while SVMs have been shown to be particularly effective for this task\cite{petprasit2015commerce}.

More recently, deep learning techniques utilizing Convolutional Neural Networks (CNNs) and Recurrent Neural Networks (RNNs) have been applied to web page classification tasks, showing improved performance over traditional machine learning algorithms.
Transformer-based models such as BERT (Bidirectional Encoder Representations from Transformers) have also been applied to web page classification tasks, showing strong performance due to their ability to capture contextual information in the text.
However, more simple models usually offer a better efficiency-performance trade-off, and are often preferred in settings where computational resources may be limited\cite{mahara2025cybersecurity}.

\subsubsection{Techniques without negative examples}\label{subsubsec:techniques_without_negative_examples}
As web page labelling is generally a manual and time consuming process,
in some cases, it may be difficult to have a complete labeled dataset consisting of both positive and negative classes (e.g. product pages vs. non-product pages).
It can be easier to just have a set of positive examples (e.g. product pages) without having a complete set of negative examples (e.g. non-product pages).
In such cases, the Positive Example Based Learning (PEBL) framework\cite{yu2004pebl} can be to  identify strong negatives from the unlabeled data and then train a classifier using the positive examples and the identified strong negatives.
This approach can be effective in situations where it is difficult to obtain a complete labeled dataset.

\subsection{Effect of noise removal}\label{subsec:effect_of_noise_removal}
In the context of web page classification, noise refers to the irrelevant information on a web page that can interfere with the accuracy of the classification process.
This can include boilerplate content, advertisements, navigation menus, and other non-content elements.
The presence of noise can lead to misclassification of web pages, as the classifier may be influenced by the irrelevant information rather than the relevant content.
By removing noise from web pages, the classifier can focus on the relevant content, which can lead to more accurate classification results.
In~\cite{shen2004web}, it was indicated that performance of web page classification improved by 10\% to 15\% after noise removal.


% effect of summarization on webpage classification
% keyword based techniques
% supervised learning techniques
% techniques without negative examples
% MarkupLM
% e-commerce web page corpus (web page classification in the context of e-commerce)